\section{Overview}\label{sec:overview}

\subsection{Aim}
AutoMarkBloom is the non-genomic sister of Genobloom. It grows flowers, starting
with peonies, from purely stochastic developmental rules with no genotype and no
pathways. The interest is in the construction method itself, so that it can
later be compared with the genetic branch to see how much realism comes from
genetics and how much from good stochastic rules.

\subsection{Two stages, two kinds of randomness}
The model follows the design principle \emph{structure first, dynamics second}.
Randomness enters at three tiers, detailed in \cref{sec:tiers}.

\begin{enumerate}[leftmargin=*]
\item \textbf{Bud formation} draws its traits per \emph{plant} but runs per
  \emph{flower}. A stochastic cellular automaton decides petal count, placement
  and outline; a second stochastic automaton on each petal decides where it
  grows; and a smooth energy packs the growing, thick petals into a bud that
  swells against its sepals. Flowers on one plant share a cultivar but are not
  clones.
\item \textbf{Flowering} opens the bud with a per-flower Markov chain over
  named channels, starting from the latent quantities the bud hands over:
  each petal grows, splays and uncurls toward a pose, and a relaxation keeps
  the petals from passing through each other.
\item \textbf{Colour} (future) places each petal or petal region on the
  parameter axis of the logistic map. Petals in a stable regime look the same
  in every flower; petals near a bifurcation amplify small differences.
\end{enumerate}

\begin{figure}[ht]
\centering
\begin{tikzpicture}[
  node distance=5mm and 7mm,
  box/.style={draw, rounded corners, align=center, font=\small, minimum height=10mm, text width=24mm},
  seed/.style={box, fill=gray!10},
  arr/.style={-{Stealth}, thick}]
\node[seed] (ps) {plant seed\\\texttt{u64}};
\node[box, right=of ps] (tr) {plant traits\\from species};
\node[seed, right=of tr] (fs) {flower seed\\plant + index};
\node[box, right=of fs] (ft) {flower traits\\vigour, layers};
\node[box, below=of ft] (sh) {nested\\shells};
\node[box, left=of sh] (ca) {shell\\automaton};
\node[box, left=of ca] (sht) {thick\\sheets};
\node[box, left=of sht] (dev) {development\\growth + packing};
\node[box, below=of dev] (bud) {packed bud\\+ channels};
\node[box, right=of bud, draw=planned, dashed] (mk) {Markov\\opening};
\node[box, right=of mk, draw=planned, dashed] (col) {logistic-map\\colour};
\draw[arr] (ps) -- (tr);
\draw[arr] (tr) -- (fs);
\draw[arr] (fs) -- (ft);
\draw[arr] (ft) -- (sh);
\draw[arr] (sh) -- (ca);
\draw[arr] (ca) -- (sht);
\draw[arr] (sht) -- (dev);
\draw[arr] (dev) -- (bud);
\draw[arr] (bud) -- (mk);
\draw[arr] (mk) -- (col);
\end{tikzpicture}
\caption{The pipeline. Solid boxes are built; dashed boxes are designed only.}
\label{fig:pipeline}
\end{figure}

\subsection{Three tiers of variation}\label{sec:tiers}
\stagebox{\statusbuilt}{\src{core/src/traits.rs}, \src{core/src/species.rs}, \src{core/src/latents.rs}}

Identical buds on every flower of a plant go too far: real flowers on one
plant share a genotype but still differ in petal count, size and placement,
through developmental noise and conditions such as temperature, and side buds
are smaller than the main one. Variation is therefore split into three tiers.

\begin{center}\small
\begin{tabularx}{\linewidth}{@{}l X@{}}
\toprule
Tier & What it sets \\
\midrule
Species & The rules and ranges: petal type (flat, cupped, rolled), spiral or
  ring arrangement, twist direction, the shape of the growth program, opening
  order, and the allowed range of every plant trait. \\
Plant & Heritable trait \emph{means}, like a cultivar: inhibition and growth
  parameters (so typical petal count and width), layer count, how double the
  flower is, thickness, bud shape, size potential, and per-layer cup and fold
  targets. \\
Flower & An overall vigour or size factor (main bud versus side bud,
  weather); its own automaton run from the plant's traits, so exact petal count
  and placement differ, with layer count varying by about one; per-petal latent
  quantities drawn around the plant's targets; then the chain's noise during
  opening. \\
\bottomrule
\end{tabularx}
\end{center}

In the code, traits are drawn per plant, and everything a flower grows (its
automaton run, growth, packing and latent quantities) is seeded per flower
(\cref{sec:seeds}). Siblings look like the same cultivar without being clones.

\paragraph{The bud hands over latent quantities.} Besides its packed geometry,
each petal leaves the bud with latent quantities, stored as named channels
(\cref{sec:channels}), that become the inputs to the opening stage
(\cref{sec:flowering}):
\begin{itemize}[leftmargin=*]
\item \emph{growth potential}: target length, width and edge growth, largely
  read from the automaton (a petal's patch width and the rows it reached stand
  in for primordium size, which sets final organ size in real flowers);
\item \emph{shape tendencies}: target fold at the base and target cupping,
  set mostly by layer;
\item \emph{timing}: when the petal starts opening, outer layers first;
\item \emph{type}: flat, cupped or rolled, and twist direction; these are
  fixed and never touched by the chain.
\end{itemize}
More latent quantities, such as colour gradients, can be added later without
changing the opening machinery.

\subsection{Why shells and sheets, not voxels}
An early idea was to grow petals as voxel sheets inside an egg-shaped volume.
Petals are thin surfaces, so a voxel grid would spend almost all of its cells on
empty space, and packing tens or hundreds of thin petals would need a very fine
grid. The model instead uses:
\begin{itemize}[leftmargin=*]
\item \emph{nested shells}: one thin egg-shaped surface per petal layer, on
  which a 2D automaton runs cheaply (\cref{sec:shells,sec:ca});
\item \emph{thick spline sheets}: each petal is a smooth mid-surface plus a
  capped thickness, so cost scales with the number of petals, not with grid
  resolution (\cref{sec:sheets});
\item \emph{a differentiable energy}: every packing cost is smooth, so exact
  gradients drive the petals into a tight, folded fit (\cref{sec:packing}).
  Differentiability also means rule settings could later be fitted to photos
  or scans of real flowers.
\end{itemize}
Discrete decisions (how many petals, which layer) stay in the automaton; the
continuous geometry is all differentiable.

\subsection{Code layout}
\begin{center}\small
\begin{tabular}{@{}ll@{}}
\toprule
Path & Contents \\
\midrule
\src{core/} & Rust core: seeds, shells, automaton, sheets, packing, JSON export \\
\src{wasm/} & Browser build of the core \\
\src{checks/} & Seed determinism tests and golden digests \\
\src{viewer/} & TypeScript and three.js viewer, glTF export \\
\src{prototype-python/} & The original 2D Python prototype \\
\src{docs/} & This document \\
\bottomrule
\end{tabular}
\end{center}
